\documentclass[11pt]{article}

\usepackage{acl}
\usepackage{times}
\usepackage{latexsym}
\usepackage[T1]{fontenc}
\usepackage[utf8]{inputenc}
\usepackage{microtype}
\usepackage{inconsolata}
\usepackage{graphicx}
\usepackage{booktabs}
\usepackage{amsmath}
\usepackage{url}

\title{Real-Time Squat Form Assessment Using MediaPipe and Temporal Conv1D Modeling}

\author{Workout Posture AI Team \\
  Open-Source Project \\
  \texttt{workout-posture-ai@project.local}}

\begin{document}
\maketitle

\begin{abstract}
This paper presents Workout Posture AI, a real-time squat form assessment system that combines pose estimation, a temporal Conv1D classifier, and rule-based biomechanical checks. The system targets six common squat errors (forward lean, heel rise, knee cave, knees over toes, shallow depth, and uneven weight) and provides immediate visual feedback, rep counting, and quality scoring. We describe the repository pipeline: landmark smoothing, body-size-invariant feature extraction, sliding-window inference, the training workflow, and the real-time user interface. The resulting system operates on live camera feeds or prerecorded videos and delivers explainable feedback without external sensors.
\end{abstract}

\section{Introduction}
Proper squat technique is fundamental for performance and injury prevention, yet deviations such as forward trunk lean and dynamic knee valgus can increase joint loading and reduce training efficacy. Recent analyses of dynamic knee valgus during squat-based tasks highlight muscle activation patterns and alignment risks that inform coaching feedback \cite{khou2024valgus}. In parallel, markerless motion capture reviews show growing clinical adoption of camera-based movement assessment \cite{lam2023mmc}. Workout Posture AI addresses this gap with a real-time, camera-only posture assessment tool that detects form errors, counts repetitions, and scores rep quality for immediate coaching.

The project integrates real-time pose estimation with a temporal Conv1D classifier and overlays rule-based checks grounded in recent biomechanics literature. The contributions of this work are: (1) a complete end-to-end pipeline for squat analysis implemented in Python; (2) a body-size-invariant feature representation with 15 interpretable features; (3) a hybrid inference approach combining deterministic rules and a learned temporal model; and (4) a real-time user interface with rep counting, quality scoring, and visual joint feedback.

\section{Related Work}
\textbf{Human pose estimation.} Recent pose estimation models and surveys demonstrate strong accuracy and practical deployment for real-time 2D keypoint tracking \cite{xu2022vitpose, zhou2023survey}. Workout Posture AI builds on these advances by using a lightweight pose pipeline to provide high-frequency body tracking for exercise feedback.

\textbf{Temporal modeling for movement classification.} Our approach leverages a compact Conv1D network over short windows of pose features to classify squat form in real time while maintaining low latency.

\textbf{Squat biomechanics and technique analysis.} Recent work on dynamic knee valgus during squat tasks motivates rule-based thresholds used in our system (e.g., forward lean, shallow depth, knee valgus), complementing learned classification for robust coaching \cite{khou2024valgus}.

\section{System Overview}
Workout Posture AI is organized into a modular pipeline. Core modules include configuration thresholds, feature extraction, a rule engine, the Conv1D model, the training pipeline, and the real-time analyzer/UI. The system accepts video from a webcam, IP camera, or file, performs pose estimation, computes features, evaluates rules and the ML model, and renders feedback overlays.

\section{Pose Acquisition and Normalization}
Each frame is processed by a real-time pose estimator, yielding 33 normalized landmarks with visibility scores. Landmarks are converted to a list of dictionaries containing \texttt{x}, \texttt{y}, \texttt{z}, and \texttt{visibility}. Frames with insufficient visibility across required joints (shoulders, hips, knees, ankles) are rejected to avoid unstable feedback. Exponential moving average smoothing is applied to both landmark coordinates and visibility scores with a configurable \(\alpha=0.6\). The smoothing factor is defined in \texttt{config.py} to provide moderate jitter reduction while maintaining responsive feedback, and it can be adjusted for different camera conditions.

Body dimensions are estimated from normalized coordinates: body height as nose-to-ankle distance and body width as shoulder-to-shoulder distance, with fallbacks for partial visibility. These measurements normalize distances and ratios to achieve body-size invariance across users.

\section{Feature Engineering}
A 15-dimensional feature vector is extracted per frame and shared between training and inference. Features capture joint angles, relative distances, and symmetry measures normalized by body height or width. This representation balances interpretability with predictive power and is described in Appendix~A.

\section{Rule-Based Form Error Detection}
Six biomechanically motivated rule checks provide immediate and explainable feedback:
\begin{itemize}
  \item \textbf{Knee cave:} detects knees moving inward relative to ankles, normalized by body width.
  \item \textbf{Forward lean:} measures trunk angle from vertical, flagging excessive flexion.
  \item \textbf{Heel rise:} compares heel height to a standing baseline to detect lifting.
  \item \textbf{Knees over toes:} uses depth (\(z\)) offsets to identify excessive forward knee travel.
  \item \textbf{Shallow depth:} evaluates hip depth ratio at the lowest squat point.
  \item \textbf{Uneven weight:} checks hip height asymmetry and knee angle imbalance.
\end{itemize}

Thresholds are derived from normalized measurements configured in \texttt{config.py}. The rules are executed only during the squat phase, and persistent violations are aggregated to score rep quality. Appendix~B lists the exact thresholds used in the implementation.

\section{Temporal Conv1D Model}
The learned classifier operates on sliding windows (15 frames \(\times\) 15 features per frame). The model stacks three Conv1D blocks (64, 128, 64 filters) with batch normalization, max pooling, and global average pooling. A dense layer with dropout (0.4) precedes the softmax output over seven classes (correct + six error classes). This architecture captures short-term temporal dynamics (e.g., knee collapse trends) while maintaining low latency for real-time feedback.

\section{Training Pipeline}
Training data are collected in \texttt{SQUAT\_VIDEOS/} and labeled via a filename-to-class mapping. The pipeline scans videos, extracts pose features, builds overlapping windows, subsamples frames for efficiency, and applies horizontal mirroring for augmentation. The model is trained with categorical cross-entropy, Adam optimization, early stopping, and learning-rate reduction; artifacts are saved to \texttt{models/} for reuse in inference.

\section{Real-Time Analysis and User Feedback}
Real-time processing is implemented by the \texttt{SquatAnalyzer} class. A finite-state machine tracks squat phases using the hip-knee depth ratio, while rule violations accumulated across frames determine rep quality (good/fair/poor). The UI overlays a skeleton with color-coded joints, highlights current issues, displays rep counts, and optionally shows ML predictions with confidence scores.

\section{Evaluation}
The repository ships with a pretrained model reporting 97.4\% validation accuracy, computed from the training pipeline's 20\% validation split and reported in the project README for the bundled model. This provides a strong baseline for seven-class discrimination on the available training videos. The repository does not publish dataset size or per-class metrics, so broader validation and more detailed reporting remain future work.

\section{Discussion and Limitations}
The hybrid design combines interpretability with statistical learning, which is advantageous for safety-critical coaching where users need actionable feedback. However, limitations remain. The system relies on 2D/monocular depth estimates, which may be less stable than multi-camera or depth sensors. Occlusions or partial visibility can reduce reliability, requiring explicit visibility checks. The dataset used for training is limited to curated squat videos and may not cover broader variations in camera angle, body type, or equipment. These constraints motivate future expansion of training data, cross-subject validation, multi-environment testing, motion-capture comparisons, sensor fusion, and adaptive personalization.

\section{Conclusion}
Workout Posture AI provides a practical, open-source pipeline for real-time squat assessment using camera-only inputs. By combining MediaPipe pose tracking, body-size-invariant features, rule-based biomechanics checks, and a temporal Conv1D classifier, the system delivers immediate, explainable coaching with minimal hardware requirements. The modular implementation supports continued research on exercise form assessment and can serve as a foundation for more comprehensive movement analysis systems.

\bibliographystyle{acl_natbib}
\bibliography{references}

\appendix
\section{Appendix A: Feature Definitions}
Table~\ref{tab:features} lists the 15 features computed per frame.

\begin{table*}[t]
  \centering
  \small
  \begin{tabular}{ll}
    \toprule
    \textbf{Feature} & \textbf{Definition (normalized)} \\
    \midrule
    left\_knee\_angle & Hip--knee--ankle angle / 180 \\
    right\_knee\_angle & Hip--knee--ankle angle / 180 \\
    back\_angle & Shoulder--hip--knee angle / 180 \\
    left\_knee\_cave\_ratio & (ankle\_x - knee\_x) / body width \\
    right\_knee\_cave\_ratio & (ankle\_x - knee\_x) / body width \\
    hip\_depth\_ratio & (avg ankle\_y - hip\_y) / body height \\
    left\_heel\_delta & (heel\_y - baseline\_heel\_y) / body height \\
    right\_heel\_delta & (heel\_y - baseline\_heel\_y) / body height \\
    shoulder\_z\_diff & |left\_shoulder\_z - right\_shoulder\_z| \\
    torso\_lean\_ratio & |shoulder\_x - hip\_x| / body width \\
    hip\_symmetry & |left\_hip\_y - right\_hip\_y| / body height \\
    knee\_angle\_diff & left\_knee\_angle - right\_knee\_angle \\
    trunk\_vertical\_angle & trunk angle from vertical / 90 \\
    knee\_forward\_ratio & avg knee\_z - avg ankle\_z \\
    hip\_lateral\_shift & (hip\_x - ankle\_x) / body width \\
    \bottomrule
  \end{tabular}
  \caption{Body-size-invariant features used for training and inference.}
  \label{tab:features}
\end{table*}

\section{Appendix B: Rule Thresholds and Rep Parameters}
Table~\ref{tab:thresholds} summarizes the thresholds used by the rule engine and rep state machine.

\begin{table}[h]
  \centering
  \small
  \begin{tabular}{ll}
    \toprule
    \textbf{Parameter} & \textbf{Value} \\
    \midrule
    Knee cave threshold & 0.12 (fraction of body width) \\
    Forward lean threshold & 50\textdegree{} from vertical \\
    Heel rise threshold & 0.04 (fraction of body height) \\
    Knees over toes threshold & 0.15 (depth offset) \\
    Shallow depth threshold & 0.18 (hip depth ratio) \\
    Uneven hip threshold & 0.06 (fraction of body height) \\
    Uneven knee threshold & 25\textdegree{} angle difference \\
    Standing hip-knee ratio & 0.23 \\
    Squat hip-knee ratio & 0.15 \\
    \bottomrule
  \end{tabular}
  \caption{Rule engine thresholds and rep detection parameters.}
  \label{tab:thresholds}
\end{table}

\section{Appendix C: Training Hyperparameters}
The default training configuration is: 100 epochs, batch size 32, validation split 0.2, early stopping patience 15, learning rate 0.001, and subsample rate 2. These parameters are configurable in \texttt{config.py} and were selected to balance training stability with real-time inference constraints.

\end{document}
